% PFC/E1/main.tex
% TA-PFC-E1: Proyecto Fin de Curso, entrega 1
% Compilacion: pdflatex main -> bibtex main -> pdflatex main -> pdflatex main

\documentclass[11pt,a4paper]{article}

% PAQUETES

\usepackage[utf8]{inputenc}
\usepackage[T1]{fontenc}
\usepackage[spanish,es-tabla]{babel}
\usepackage[margin=2.5cm]{geometry}
\usepackage{enumitem}
\usepackage{booktabs}
\usepackage{longtable}
\usepackage{tabularx}
\usepackage{array}
\usepackage{graphicx}
\usepackage{xcolor}
\usepackage{float}
\usepackage{url}
\usepackage{seqsplit}
\usepackage[hidelinks]{hyperref}
\usepackage{setspace}
\usepackage{fancyhdr}

\setstretch{1.1}
\setlength{\parskip}{4pt}

% CONFIGURACIONES Y COMANDOS

% Marcador de contenido por completar (eliminar todos antes de entregar)
\newcommand{\pendiente}[1]{\textcolor{red}{[#1]}}

% Figura con respaldo: si el archivo no existe, muestra un recuadro y compila igual
\newcommand{\figuraopcional}[3]{%
  \IfFileExists{figuras/#1}
  {\includegraphics[width=#2]{figuras/#1}}
  {\fbox{\parbox{#2}{\centering\vspace{2cm}
  \pendiente{#3}\vspace{2cm}}}}%
}

\newcolumntype{Y}{>{\raggedright\arraybackslash}X}

\pagestyle{fancy}
\setlength{\headheight}{14pt}
\fancyhf{}
\lhead{UTEQ \textperiodcentered{} Ingeniería de Requerimientos (ISR-401)}
\rhead{TA-PFC-E1}
\cfoot{\thepage}

\begin{document}

% PORTADA

\begin{titlepage}
    \centering

    {\large \textbf{UNIVERSIDAD TÉCNICA ESTATAL DE QUEVEDO}}\\[0.3cm]
    {\large Facultad de Ciencias de la Computación}\\
    {\large Carrera de Software}\\[1.5cm]

    {\Large \textbf{INGENIERÍA DE REQUERIMIENTOS}}\\[0.5cm]
    {\large ISR-401}\\[1.5cm]

    {\LARGE \textbf{TA-PFC-E1}}\\[0.9cm]

    \rule{0.85\textwidth}{0.5pt}\\[0.5cm]

    {\large \textbf{Proyecto Fin de Curso, entrega 1: visión, alcance, partes interesadas, contexto y plan de elicitación}}\\[0.3cm]

    {\normalsize Grupo: \textbf{HLRS}}\\

    \rule{0.85\textwidth}{0.5pt}\\[1.5cm]

    \begin{flushleft}
        \textbf{Integrantes:}
        \begin{itemize}[noitemsep,topsep=2pt]
            \item Herrera Moran Gilmar Javier
            \item Llerena Abril Angelina Julexy
            \item Rivera Rodriguez Joseph Jeremy
            \item Saltos Tello Joseph Orlando
        \end{itemize}

        \vspace{0.2cm}

        \textbf{Nivel:} Cuarto\\[0.4cm]
        \textbf{Paralelo:} A\\[0.4cm]
        \textbf{Docente:} Gleiston Cicerón Guerrero Ulloa, PhD\\[0.4cm]
        \textbf{Modalidad:} Grupal\\[0.4cm]
        \textbf{Fecha límite de entrega:} sábado 21 de noviembre de 2026, 23:59
    \end{flushleft}

    \vfill

    {\large Quevedo -- Ecuador}\\
    {\large 2026}

\end{titlepage}

% INDICES

\tableofcontents
\listoffigures
\listoftables
\newpage

% 1. INTRODUCCION Y ORGANIZACION

\section{Introducción}

El presente documento corresponde a la primera entrega del Proyecto Fin de Curso de la asignatura Ingeniería de Requerimientos (ISR-401). Su propósito es establecer la visión inicial, el alcance, las partes interesadas, el contexto y el plan de elicitación de una propuesta de software orientada a apoyar la gestión de un local de manualidades.

La propuesta consiste en desarrollar una aplicación web que permita organizar información relacionada con clientes, proveedores, empleados, productos terminados, pedidos, ventas, pagos, abonos y entregas. La iniciativa parte de la necesidad de analizar los procesos de gestión del establecimiento y determinar de qué manera un sistema informático podría contribuir a su organización y seguimiento.

Esta entrega establece una planificación preliminar de las actividades de ingeniería de requisitos. Por tanto, los procesos descritos representan el ámbito que se pretende estudiar y deberán contrastarse con la información proporcionada por los responsables y participantes de la organización.

La documentación se estructura considerando los procesos de ingeniería de requisitos y la identificación de las partes interesadas, el contexto y las necesidades del sistema, de acuerdo con los lineamientos de la norma ISO/IEC/IEEE 29148:2018 \cite{iso29148}.


\section{Organización real}

\subsection{Descripción de la organización}

Pauly Creative's Craft es un emprendimiento dedicado a la elaboración y comercialización de productos de manualidades, ubicado en la ciudad de Quevedo, provincia de Los Ríos, Ecuador. Su actividad comercial se relaciona con el ámbito creativo y artesanal, ofreciendo productos asociados a las manualidades.

La organización constituye el caso de estudio seleccionado para el desarrollo del Proyecto Fin de Curso, cuyo propósito es analizar las necesidades relacionadas con la gestión comercial y administrativa del emprendimiento, a fin de establecer los requisitos de una propuesta de software.

La información presentada en este apartado se limita a datos generales de identificación y descripción de la organización. Los detalles relacionados con su funcionamiento interno, información comercial específica, datos personales y demás información confidencial no se exponen públicamente y serán tratados únicamente en el contexto académico autorizado.

\textbf{Datos de identificación de la organización:}

\begin{itemize}
	\item \textbf{Nombre comercial:} Pauly Creative's Craft.
	\item \textbf{Actividad económica:} Elaboración y comercialización de productos de manualidades.
	\item \textbf{Ubicación:} Quevedo, provincia de Los Ríos, Ecuador.
	\item \textbf{Responsable:} Propietaria del emprendimiento.
	\item \textbf{Tipo de organización:} Emprendimiento comercial dedicado al ámbito de las manualidades.
\end{itemize}

Para el desarrollo del proyecto se utilizará información obtenida mediante las actividades de elicitación autorizadas por la organización. Los datos que puedan identificar a clientes, proveedores u otras personas, así como información interna sensible, serán excluidos de los documentos y repositorios públicos.

\subsection{Problema identificado}

El proyecto contempla el análisis de las necesidades relacionadas con la gestión comercial y administrativa de un emprendimiento dedicado a la elaboración y comercialización de productos de manualidades.

Como punto de partida, se plantea estudiar aspectos relacionados con la gestión de pedidos, clientes, productos, ventas y demás actividades vinculadas al contexto del establecimiento, con el propósito de identificar posibles necesidades de información y oportunidades de mejora.

La problemática específica no se considera confirmada en esta etapa, debido a que deberá determinarse mediante las técnicas de elicitación planificadas y la información que la organización autorice proporcionar.

Por tanto, cualquier necesidad, dificultad o deficiencia identificada será documentada a partir de información verificable obtenida durante el levantamiento de requisitos, evitando divulgar datos internos o confidenciales del emprendimiento.

\subsection{Acuerdo con la organización}

La organización seleccionada como caso de estudio es Pauly Creative's Craft. La participación se coordinará con su propietaria, a quien se comunicará el propósito académico del proyecto, las actividades previstas y el tratamiento que se dará a la información obtenida.

La aceptación de la organización se respaldará mediante la carta de aceptación correspondiente. Este documento será resguardado por el grupo de trabajo y su publicación estará sujeta a las condiciones de autorización establecidas por la organización.

El registro del primer contacto y la identificación de las evidencias se mantendrán en el archivo \path{evidencias/MANIFIESTO.md}. La documentación pública del proyecto contendrá únicamente información autorizada y no incluirá datos personales ni detalles confidenciales del establecimiento.

\textbf{Estado del acuerdo:} Primer contacto realizado el 19/11/2026. La formalización de la participación se respaldará mediante la carta de aceptación correspondiente.
```
% Visión y alcance
\section{Visión y alcance}

\subsection{Declaración de visión}

Desarrollar una aplicación web que facilite la gestión de un local de manualidades, reuniendo la información de clientes, proveedores, productos terminados, pedidos, ventas, pagos, abonos y entregas. Se dará mayor importancia al registro y seguimiento de los pedidos personalizados, desde que el cliente los solicita hasta que se entregan.

La aplicación busca facilitar el registro, la consulta y el seguimiento de las actividades del establecimiento. De esta manera, se contará con información organizada que ayude a tomar decisiones y a llevar un mejor control de las actividades comerciales.

El sistema estará dirigido al propietario del establecimiento como único usuario. Los clientes no tendrán acceso al sistema, pero participarán proporcionando información durante el levantamiento de requisitos.

\subsection{Objetivos de negocio}

Los objetivos de negocio se presentan en la Tabla~\ref{tab:objetivos}. Las metas propuestas son valores iniciales que servirán como referencia para el diseño del sistema y deberán revisarse con el responsable del establecimiento antes de establecerse como metas definitivas.

\begin{table}[H]
\centering
\caption{Objetivos de negocio medibles.}
\label{tab:objetivos}
\small
\begin{tabularx}{\textwidth}{@{}lYYYY@{}}
\toprule
ID & Objetivo & Indicador & Meta propuesta & Fuente de verificación \\
\midrule
OBJ-01 &
Reunir en el sistema la información de clientes, proveedores y productos terminados. &
Porcentaje de registros gestionados mediante el sistema. &
Al menos 90\% de los registros considerados dentro del alcance. &
Pruebas funcionales y revisión de registros. \\

OBJ-02 &
Facilitar el registro y seguimiento de los pedidos personalizados desde que se reciben hasta que se entregan. &
Porcentaje de pedidos personalizados con estado actualizado y consultable. &
Al menos 90\% de los pedidos registrados con seguimiento. &
Pruebas de gestión de pedidos. \\

OBJ-03 &
Organizar la información de ventas, pagos, abonos y entregas para facilitar su consulta. &
Porcentaje de operaciones registradas y consultables. &
Al menos 90\% de las operaciones incluidas en el sistema. &
Pruebas funcionales y reportes generados. \\
\bottomrule
\end{tabularx}
\end{table}

\subsection{Alcance}

\subsubsection{Incluido en el sistema}

El sistema será una aplicación web que ayudará a gestionar las actividades del establecimiento. En su primera versión contará con los siguientes módulos:

\begin{itemize}
    \item \textbf{Gestión de clientes:} permitirá registrar, consultar y actualizar la información necesaria para atender a los clientes.
    \item \textbf{Gestión de proveedores:} permitirá registrar y consultar la información de los proveedores.
    \item \textbf{Gestión de usuarios:} permitirá el ingreso del propietario al sistema como único usuario administrador.
    \item \textbf{Gestión de productos:} permitirá registrar y actualizar los productos terminados, así como consultar su disponibilidad.
    \item \textbf{Gestión de pedidos personalizados:} permitirá registrar las características solicitadas por el cliente, el estado del pedido, los valores pagados y pendientes, y la fecha prevista de entrega.
    \item \textbf{Gestión de ventas:} permitirá registrar y consultar las ventas de productos terminados y de pedidos personalizados que ya hayan sido entregados.
    \item \textbf{Gestión de pagos y abonos:} permitirá registrar los pagos realizados y los valores que aún estén pendientes.
    \item \textbf{Gestión de entregas:} permitirá registrar y llevar un seguimiento de las entregas de los pedidos.
    \item \textbf{Reportes:} permitirá generar resúmenes de información sobre pedidos, ventas, pagos y entregas.
\end{itemize}


\subsubsection{Excluido del sistema}

Para esta primera etapa, no se incluirán las siguientes funcionalidades:

\begin{itemize}
    \item Control de inventario de materiales e insumos; únicamente se manejará el inventario de productos terminados.
    \item Acceso de los clientes al sistema para realizar o consultar pedidos.
    \item Integración con entidades bancarias y plataformas de pago externas.
    \item Procesamiento automático de transacciones bancarias.
    \item Gestión contable y tributaria completa del establecimiento.
    \item Administración de nómina y cálculo de remuneraciones.
    \item Comercio electrónico público con carrito de compras y pasarela de pago.
    \item Aplicación móvil nativa independiente.
    \item Control automatizado de maquinaria o procesos físicos de fabricación.
    \item Integración con servicios externos de transporte o mensajería.
\end{itemize}

Estas exclusiones podrán revisarse si durante el análisis de las necesidades se determina que alguna funcionalidad es necesaria y se justifica su incorporación.

\subsection{Restricciones conocidas}

Para el desarrollo del proyecto se han considerado las siguientes restricciones:

\begin{itemize}
    \item \textbf{Tiempo:} el desarrollo se realizará dentro del período académico establecido para la asignatura.
    \item \textbf{Disponibilidad:} las actividades de levantamiento de requisitos dependerán de los horarios y del tiempo disponible del propietario y de los demás participantes.
    \item \textbf{Recursos:} la propuesta deberá adaptarse a los recursos técnicos y humanos disponibles para el desarrollo del proyecto académico.
    \item \textbf{Acceso a información:} no se solicitarán contraseñas, datos bancarios ni listados completos de clientes o proveedores.
    \item \textbf{Privacidad:} la información personal y confidencial se manejará respetando los principios de protección de datos aplicables.
    \item \textbf{Conectividad:} al ser una aplicación web, será necesario contar con conexión a internet para utilizarla.
    \item \textbf{Validación:} las funcionalidades y sus prioridades deberán revisarse con el propietario y los demás participantes del proyecto.
\end{itemize}


% 3. PARTES INTERESADAS

\section{Partes interesadas}

\subsection{Identificación}

Se identifican al menos ocho partes interesadas, expresadas como roles y no como nombres (Tabla~\ref{tab:interesados}). La norma ISO/IEC/IEEE 29148:2018~\cite{iso29148} exige identificar a las partes interesadas y definir el límite del sistema antes de definir los requisitos.

\begin{table}[H]
\centering
\caption{Partes interesadas, interés, influencia y disponibilidad.}
\label{tab:interesados}
\small
\begin{tabularx}{\textwidth}{@{}lYYYY@{}}
\toprule
ID & Rol & Interés & Poder o influencia & Disponibilidad \\
\midrule
PI-01 & \pendiente{...} & \pendiente{...} & \pendiente{...} & \pendiente{...} \\
PI-02 & \pendiente{...} & \pendiente{...} & \pendiente{...} & \pendiente{...} \\
PI-03 & \pendiente{...} & \pendiente{...} & \pendiente{...} & \pendiente{...} \\
PI-04 & \pendiente{...} & \pendiente{...} & \pendiente{...} & \pendiente{...} \\
PI-05 & \pendiente{...} & \pendiente{...} & \pendiente{...} & \pendiente{...} \\
PI-06 & \pendiente{...} & \pendiente{...} & \pendiente{...} & \pendiente{...} \\
PI-07 & \pendiente{...} & \pendiente{...} & \pendiente{...} & \pendiente{...} \\
PI-08 & \pendiente{...} & \pendiente{...} & \pendiente{...} & \pendiente{...} \\
\bottomrule
\end{tabularx}
\end{table}

\subsection{Matriz poder--interés}

\begin{figure}[H]
\centering
\figuraopcional{matriz_poder_interes.pdf}{0.85\textwidth}{Exportar la matriz poder--interés desde modelos/}
\caption{Matriz poder--interés de las partes interesadas.}
\label{fig:poder-interes}
\end{figure}

La Figura~\ref{fig:poder-interes} ubica a cada parte interesada en su cuadrante.

\pendiente{Interpretación de la matriz.}

\subsection{Estrategia de participación}

La Tabla~\ref{tab:estrategia} define la estrategia de participación de cada parte interesada según su cuadrante en la matriz.

\begin{table}[H]
\centering
\caption{Estrategia de participación por parte interesada.}
\label{tab:estrategia}
\small
\begin{tabularx}{\textwidth}{@{}lYYY@{}}
\toprule
ID & Cuadrante & Estrategia de participación & Frecuencia y canal \\
\midrule
PI-01 & \pendiente{...} & \pendiente{...} & \pendiente{...} \\
PI-02 & \pendiente{...} & \pendiente{...} & \pendiente{...} \\
PI-03 & \pendiente{...} & \pendiente{...} & \pendiente{...} \\
PI-04 & \pendiente{...} & \pendiente{...} & \pendiente{...} \\
PI-05 & \pendiente{...} & \pendiente{...} & \pendiente{...} \\
PI-06 & \pendiente{...} & \pendiente{...} & \pendiente{...} \\
PI-07 & \pendiente{...} & \pendiente{...} & \pendiente{...} \\
PI-08 & \pendiente{...} & \pendiente{...} & \pendiente{...} \\
\bottomrule
\end{tabularx}
\end{table}

% 4. CONTEXTO

\section{Contexto del sistema}

\subsection{Límite del sistema y diagrama de contexto}

\begin{figure}[H]
\centering
\figuraopcional{diagrama_contexto.pdf}{0.9\textwidth}{Exportar el diagrama de contexto desde modelos/}
\caption{Diagrama de contexto del sistema.}
\label{fig:contexto}
\end{figure}

La Figura~\ref{fig:contexto} muestra el sistema, sus actores externos y los flujos de información.

\pendiente{Explicación del límite del sistema.}

\subsection{Perspectivas del sílabo}

La Tabla~\ref{tab:perspectivas} resume los hallazgos obtenidos en la organización para cada una de las ocho perspectivas del sílabo.

\begin{table}[H]
\centering
\caption{Las ocho perspectivas del sílabo con hallazgos de la organización real.}
\label{tab:perspectivas}
\small
\begin{tabularx}{\textwidth}{@{}lYYY@{}}
\toprule
N.º & Perspectiva (según sílabo) & Hallazgo en la organización & Evidencia \\
\midrule
1 & \pendiente{...} & \pendiente{...} & \pendiente{...} \\
2 & \pendiente{...} & \pendiente{...} & \pendiente{...} \\
3 & \pendiente{...} & \pendiente{...} & \pendiente{...} \\
4 & \pendiente{...} & \pendiente{...} & \pendiente{...} \\
5 & \pendiente{...} & \pendiente{...} & \pendiente{...} \\
6 & \pendiente{...} & \pendiente{...} & \pendiente{...} \\
7 & \pendiente{...} & \pendiente{...} & \pendiente{...} \\
8 & \pendiente{...} & \pendiente{...} & \pendiente{...} \\
\bottomrule
\end{tabularx}
\end{table}

% 5. PLAN DE ELICITACION

\section{Plan de elicitación}

\subsection{Selección de técnicas}

La selección de técnicas se fundamenta en la fuente y en el tipo de conocimiento que se busca obtener, según lo planteado por Pohl~\cite{pohl2025} y por el sílabo de nivel fundamental del IREB~\cite{ireb2026}. La Tabla~\ref{tab:tecnicas} la presenta por parte interesada.

\begin{table}[H]
\centering
\caption{Técnica de elicitación por parte interesada y justificación.}
\label{tab:tecnicas}
\small
\begin{tabularx}{\textwidth}{@{}lYYY@{}}
\toprule
ID & Técnica & Tipo de conocimiento / fuente & Justificación \\
\midrule
PI-01 & \pendiente{...} & \pendiente{...} & \pendiente{...} \\
PI-02 & \pendiente{...} & \pendiente{...} & \pendiente{...} \\
PI-03 & \pendiente{...} & \pendiente{...} & \pendiente{...} \\
PI-04 & \pendiente{...} & \pendiente{...} & \pendiente{...} \\
PI-05 & \pendiente{...} & \pendiente{...} & \pendiente{...} \\
PI-06 & \pendiente{...} & \pendiente{...} & \pendiente{...} \\
PI-07 & \pendiente{...} & \pendiente{...} & \pendiente{...} \\
PI-08 & \pendiente{...} & \pendiente{...} & \pendiente{...} \\
\bottomrule
\end{tabularx}
\end{table}

\subsection{Cronograma de sesiones}

Las sesiones de elicitación se programan según la Tabla~\ref{tab:sesiones}.

\begin{table}[H]
\centering
\caption{Cronograma de sesiones de elicitación.}
\label{tab:sesiones}
\small
\begin{tabularx}{\textwidth}{@{}lllYY@{}}
\toprule
Sesión & Fecha & Duración & Participantes (roles) & Técnica e instrumento \\
\midrule
S1 & \pendiente{...} & \pendiente{...} & \pendiente{...} & \pendiente{...} \\
S2 & \pendiente{...} & \pendiente{...} & \pendiente{...} & \pendiente{...} \\
S3 & \pendiente{...} & \pendiente{...} & \pendiente{...} & \pendiente{...} \\
S4 & \pendiente{...} & \pendiente{...} & \pendiente{...} & \pendiente{...} \\
\bottomrule
\end{tabularx}
\end{table}

\subsection{Instrumentos}

Los instrumentos completos se encuentran en \path{PFC/E1/instrumentos/} y se resumen a continuación.

\subsubsection{Guion de entrevista}

\pendiente{Objetivo, estructura, preguntas clave y relación con los objetivos de negocio (OBJ-01 a OBJ-03).}

\subsubsection{Cuestionario}

\pendiente{Objetivo, población, tipo de preguntas y escala.}

\subsubsection{Guía de observación o de taller}

\pendiente{Objetivo, aspectos a observar o dinámica del taller, registro de resultados.}

\subsection{Trazabilidad entre partes interesadas, técnicas e instrumentos}

La Tabla~\ref{tab:trazabilidad} relaciona cada objetivo de negocio con las partes interesadas, las técnicas y los instrumentos correspondientes.

\begin{table}[H]
\centering
\caption{Trazabilidad entre objetivos, partes interesadas, técnicas e instrumentos.}
\label{tab:trazabilidad}
\small
\begin{tabularx}{\textwidth}{@{}lYYY@{}}
\toprule
Objetivo & Partes interesadas & Técnica & Instrumento \\
\midrule
OBJ-01 & \pendiente{...} & \pendiente{...} & \pendiente{...} \\
OBJ-02 & \pendiente{...} & \pendiente{...} & \pendiente{...} \\
OBJ-03 & \pendiente{...} & \pendiente{...} & \pendiente{...} \\
\bottomrule
\end{tabularx}
\end{table}

% 6. ETICA Y DATOS

\section{Ética y protección de datos}

\subsection{Consentimiento informado}

Para las actividades de elicitación se utilizará un formulario de consentimiento informado dirigido a personas mayores de 18 años que participen voluntariamente en el proyecto. Este documento explicará el propósito académico del estudio, las actividades previstas, los posibles riesgos, los beneficios, el tratamiento de la información y el derecho de cada participante a retirarse sin consecuencias. También incluirá autorizaciones independientes para la participación, la grabación de audio y video, la toma de fotografías del entorno de trabajo, la cita de frases textuales y la utilización de información anonimizada en trabajos académicos, publicaciones y repositorios abiertos.

La autorización para participar será independiente de los permisos relacionados con registros audiovisuales, fotográficos y otros usos de la información, por lo que cada persona podrá participar sin aceptar dichas autorizaciones adicionales. El modelo de consentimiento se encuentra en el documento correspondiente del proyecto y se incorpora como referencia en el Anexo~\ref{anx:consentimiento}. El tratamiento de los datos personales considerará la Ley Orgánica de Protección de Datos Personales del Ecuador \cite{lopdp2021} y su Reglamento General \cite{reglamentolopdp2023}.

\subsection{Anonimización de datos}

\pendiente{Procedimiento: códigos de participante, eliminación de datos identificables, tratamiento de transcripciones.}

\subsection{Custodia de evidencias}

\pendiente{Carpeta institucional privada (cuenta @uteq.edu.ec) compartida con el docente, responsables, control de acceso y plazo de conservación. En el repositorio se publica solo el manifiesto con nombre de archivo, fecha, duración y huella SHA-256 (Tabla~\ref{tab:manifiesto}).}

\subsection{Principios éticos}

\pendiente{Aplicación del Código de Ética de la ACM.}

% 7. GESTION DEL PROYECTO

\section{Gestión del proyecto}

\subsection{Roles del equipo}

Los roles y responsabilidades de los integrantes se detallan en la Tabla~\ref{tab:roles}.

\begin{table}[H]
\centering
\caption{Roles y responsabilidades del equipo.}
\label{tab:roles}
\small
\begin{tabularx}{\textwidth}{@{}lYY@{}}
\toprule
Integrante & Rol & Responsabilidades \\
\midrule
Herrera Moran Gilmar Javier & Observación y evidencias & Guía de observación, matriz poder--interés, anonimización y manifiesto de evidencias \\
Llerena Abril Angelina Julexy & Entrevista y ética & Guion de entrevista, consentimiento informado, principios éticos y referencias \\
Rivera Rodriguez Joseph Jeremy & Cuestionario y modelado & Cuestionario, diagrama de contexto y README \\
Saltos Tello Joseph Orlando & Coordinación & Contacto con el local, carta de aceptación, PDF, entrega en el SGA y declaración de uso de IA \\
\bottomrule
\end{tabularx}
\end{table}

\subsection{Cronograma hasta la semana 17}

La Tabla~\ref{tab:cronograma} presenta las actividades, productos y responsables del PFC hasta la semana 17.

\begin{table}[H]
\centering
\caption{Cronograma del PFC hasta la semana 17.}
\label{tab:cronograma}
\small
\begin{tabularx}{\textwidth}{@{}lYYY@{}}
\toprule
Semana & Actividad & Producto & Responsable \\
\midrule
4 & Entrega 1 (esta entrega) & \pendiente{...} & \pendiente{...} \\
5 & \pendiente{...} & \pendiente{...} & \pendiente{...} \\
6 & \pendiente{...} & \pendiente{...} & \pendiente{...} \\
7 & \pendiente{...} & \pendiente{...} & \pendiente{...} \\
8 & \pendiente{...} & \pendiente{...} & \pendiente{...} \\
9 & \pendiente{...} & \pendiente{...} & \pendiente{...} \\
10 & \pendiente{...} & \pendiente{...} & \pendiente{...} \\
11 & \pendiente{...} & \pendiente{...} & \pendiente{...} \\
12 & \pendiente{...} & \pendiente{...} & \pendiente{...} \\
13 & \pendiente{...} & \pendiente{...} & \pendiente{...} \\
14 & \pendiente{...} & \pendiente{...} & \pendiente{...} \\
15 & \pendiente{...} & \pendiente{...} & \pendiente{...} \\
16 & \pendiente{...} & \pendiente{...} & \pendiente{...} \\
17 & \pendiente{...} & \pendiente{...} & \pendiente{...} \\
\bottomrule
\end{tabularx}
\end{table}

\subsection{Plan de comunicación con la organización}

La comunicación con la organización se planifica según la Tabla~\ref{tab:comunicacion}.

\begin{table}[H]
\centering
\caption{Plan de comunicación con la organización.}
\label{tab:comunicacion}
\small
\begin{tabularx}{\textwidth}{@{}lYYYY@{}}
\toprule
Información & Destinatario (rol) & Canal & Frecuencia & Responsable \\
\midrule
\pendiente{...} & \pendiente{...} & \pendiente{...} & \pendiente{...} & \pendiente{...} \\
\pendiente{...} & \pendiente{...} & \pendiente{...} & \pendiente{...} & \pendiente{...} \\
\bottomrule
\end{tabularx}
\end{table}

\subsection{Registro de riesgos}

Los riesgos se formulan de manera específica para la organización y el proyecto, siguiendo los principios de gestión de proyectos de la Guía del PMBOK, 7.ª edición~\cite{pmbok7}. La Tabla~\ref{tab:riesgos} los registra junto con su respuesta y responsable.

\begin{table}[H]
\centering
\caption{Registro de riesgos.}
\label{tab:riesgos}
\small
\begin{tabularx}{\textwidth}{@{}lYcccYY@{}}
\toprule
ID & Riesgo & Prob. & Impacto & Nivel & Respuesta & Responsable \\
\midrule
R-01 & \pendiente{...} & \pendiente{...} & \pendiente{...} & \pendiente{...} & \pendiente{...} & \pendiente{...} \\
R-02 & \pendiente{...} & \pendiente{...} & \pendiente{...} & \pendiente{...} & \pendiente{...} & \pendiente{...} \\
R-03 & \pendiente{...} & \pendiente{...} & \pendiente{...} & \pendiente{...} & \pendiente{...} & \pendiente{...} \\
\bottomrule
\end{tabularx}
\end{table}

% 8. ANALISIS Y CONCLUSIONES

\section{Análisis crítico}

\pendiente{Interpretación de lo realizado, alternativas consideradas (por ejemplo, técnicas descartadas), limitaciones de la entrega y coherencia entre partes interesadas, técnicas e instrumentos.}

\section{Conclusiones}

\pendiente{Conclusiones derivadas de lo realizado en la organización, sin afirmaciones genéricas.}

% 9. DECLARACION DE USO DE IA

\section{Declaración de uso de inteligencia artificial}

\pendiente{Si se usó una herramienta de IA solo como apoyo (ortografía o errores de compilación), indicar: herramienta, versión, propósito y cómo se verificó lo obtenido. Si no se usó, declararlo expresamente.}


% ANEXOS

\appendix

\section{Modelo de consentimiento informado}
\label{anx:consentimiento}

\pendiente{Título del estudio, responsable y contacto institucional, finalidad, participación voluntaria, derecho a retirarse, tratamiento y plazo de conservación de datos, anonimización, derechos del titular (acceso, rectificación, eliminación), y espacio para firma y fecha, sin publicar ningún consentimiento firmado.}

\section{Registro de evidencias (manifiesto)}
\label{anx:manifiesto}

La Tabla~\ref{tab:manifiesto} registra las evidencias privadas.

\begin{table}[H]
\centering
\caption{Manifiesto de evidencias privadas.}
\label{tab:manifiesto}
\footnotesize
\begin{tabularx}{\textwidth}{@{}YYllY@{}}
\toprule
Archivo & Descripción & Fecha & Duración & SHA-256 \\
\midrule
\pendiente{...} & Carta de aceptación & \pendiente{...} & --- & \pendiente{hash SHA-256} \\
\pendiente{...} & Registro de primer contacto & \pendiente{...} & \pendiente{...} & \pendiente{hash SHA-256} \\
\bottomrule
\end{tabularx}
\end{table}

\section{Instrumentos de elicitación}

\pendiente{Referencia a los archivos de PFC/E1/instrumentos/ (guion de entrevista, cuestionario, guía de observación o taller).}

\bibliographystyle{IEEEtran}
\bibliography{referencias}

\end{document}